\begin{claim}{F.25}{Proved}{The actual affine marginal images with a common scalar field}
Retain $\mathcal K_\pm\leq R_\pm\leq V_\pm$ and their
actual maps from F.19--F.24. Let $k$ be a finite field
of characteristic $p$, and let
$\rho:V_\pm\to\operatorname{Aff}_1(k)$ be continuous.
Write $\chi$ for its scalar character and put
\[
 \begin{gathered}
 M=\rho(\mathcal K_\pm),\qquad N=\rho(R_\pm),\qquad
 D=\chi(V_\pm),\qquad C=\chi(R_\pm),\\
 U=\{b\in k:(b,1)\in M\},\qquad
 k_0=\mathbb F_p(D),\qquad k_1=\mathbb F_p(C).
 \end{gathered}                                                \tag{F.25.1}
\]
The subgroup $U$ is a vector subspace of $k$ over
$k_1$. For the class of these representations in
which $U$ is also a vector subspace over $k_0$, the
actual finite images agree:
\[
                              M=N.                            \tag{F.25.2}
\]
Consequently the equality holds for every representation
in the explicitly specified class
\[
 \left\{\rho:\ C=\{1\}\ \text{or}\ k_1=k_0\right\}.
                                                               \tag{F.25.3}
\]
In particular, for every prime $p$ and every continuous
homomorphism to the full affine group of its prime field,
\[
 \rho(\mathcal K_\pm)=\rho(R_\pm)
       \qquad\bigl(\rho:V_\pm\to\operatorname{Aff}_1(\mathbb F_p)
                       \text{ continuous}\bigr).             \tag{F.25.4}
\]
All comparison maps, the choice of an actual scalar
lift in $M$, and the coefficient maps used in the proof
are specified below. The finite scalar-field class
$k_1\subsetneq k_0$ with $U$ failing to be a $k_0$-subspace
remains outside the proved equality; full marginal
surjectivity and $K^+=D^+$ remain unresolved.
\end{claim}

\begin{proof}
Fix either sign throughout. First transport the F.21
normal-generation statement through the actual continuous
surjection $p_\pm:D^+\to R_\pm$ of F.22.6. Its restriction
to $K^+$ has image $\mathcal K_\pm$. The image of the
closed normal closure is the closed normal closure of
the image: one inclusion follows by applying the map
to conjugates, and the other by taking the inverse
image of a closed normal subgroup containing the image.
The image is closed because the source is compact and
the target Hausdorff. Thus
\[
 \overline{\langle\!\langle\mathcal K_\pm
                         \rangle\!\rangle}_{R_\pm}=R_\pm,
 \qquad
 \langle\!\langle M\rangle\!\rangle_N=N,
 \qquad \chi(\mathcal K_\pm)=C.                               \tag{F.25.5}
\]
The second equality applies the actual map
$\rho|R_\pm$; the third applies its scalar projection
to the finite abelian group $k^\times$. When $C=\{1\}$,
the group $N$ is a subgroup of the abelian translation
group, so (F.25.5) already gives $M=N$.

We now take $C\ne\{1\}$. The set $U$ is an additive
subgroup of $k$, hence an $\mathbb F_p$-vector subspace.
Conjugation of a translation by an actual element of
$M$ with scalar $c$ is given by the literal formula
\[
             (b,c)(u,1)(b,c)^{-1}=(cu,1).                    \tag{F.25.6}
\]
The last equality of (F.25.5) supplies such an element
for every $c\in C$. Consequently $U$ is stable under
every scalar in $C$, hence under the ring
$\mathbb F_p[C]$. This finite subring of $k$ is a field:
multiplication by any nonzero element is an injection
of its finite set and therefore is onto, giving its
inverse in the ring. Thus $\mathbb F_p[C]=k_1$ and
$U$ is a $k_1$-subspace, as asserted.

The finite multiplicative subgroup $C$ is cyclic.
One direct verification chooses an element of maximal
order $n$: the prime-power components of element
orders in the abelian group show that every element
has order dividing $n$, since otherwise their commuting
components would give a larger order. All elements
then satisfy $X^n-1=0$ in the field, so $|C|\leq n$.
The chosen cyclic subgroup has $n$ elements, proving
$|C|=n$. Choose a generator $c\ne1$ and an actual
lift $m=(b,c)\in M$. Define
\[
 z=\frac{b}{1-c},\qquad \tau_z(w)=w+z,
 \qquad
 \rho_z(g)=\tau_z^{-1}\rho(g)\tau_z
           =\bigl(d(g)+(\chi(g)-1)z,\chi(g)\bigr),             \tag{F.25.7}
\]
where $\rho(g)=(d(g),\chi(g))$. The maps $\tau_z$ and
$\tau_z^{-1}=\tau_{-z}$ are explicit inverse affine
bijections of $k$, and their conjugations are inverse
group automorphisms of $\operatorname{Aff}_1(k)$.
The chosen lift now equals $(0,c)$. Conjugating by a
translation fixes every element of $T_k$, so the
translation intersection remains exactly $U$.

Let $M_z=\rho_z(\mathcal K_\pm)$ and $N_z=\rho_z(R_\pm)$.
Every element of $M_z$ has scalar $c^j$; multiplication
by its element $(0,c^{-j})$ leaves a translation in $U$.
Conversely $M_z$ contains all translations in $U$ and
the generator $(0,c)$. We have therefore proved the
literal equality and the associated inverse coordinates
\[
 \begin{aligned}
 M_z&=U\rtimes C,\\
 U\rtimes C&\longrightarrow M_z,&(u,a)&\longmapsto(u,a),\\
 M_z&\longrightarrow U\rtimes C,&(b,a)&\longmapsto(b,a).
 \end{aligned}                                                \tag{F.25.8}
\]
The group law is $(u,a)(v,a')=(u+av,aa')$, with
$U$ stable under $C$ by (F.25.6). Thus the complement
used here has been constructed from the actual
element $m\in M$.

Now let $U$ be a $k_0$-subspace, as in the first
assertion. The field $k$ is a finite-dimensional vector
space over $k_0$. Let
\[
 \mathscr L_U=\{\ell\in\operatorname{Hom}_{k_0}(k,k_0):
                                      \ell(U)=0\}.
\]
These exact coefficient maps satisfy
\[
                         U=\bigcap_{\ell\in\mathscr L_U}
                                           \ker\ell.          \tag{F.25.9}
\]
For a proof, extend a $k_0$-basis of $U$ to a basis of
$k$. Coordinate projections on the additional basis
vectors belong to $\mathscr L_U$ and detect every
nonzero coordinate outside $U$; the reverse inclusion
is part of the definition.

Put $d_z(g)=d(g)+(\chi(g)-1)z$, the translation
coordinate in (F.25.7). The affine multiplication law
makes it a continuous $k$-valued crossed homomorphism.
For each $\ell\in\mathscr L_U$ the composite
\[
 \ell d_z:V_\pm\longrightarrow(k_0,+),\qquad
 (\ell d_z)(gh)=(\ell d_z)(g)
                         +\chi(g)(\ell d_z)(h)               \tag{F.25.10}
\]
is an actual continuous crossed homomorphism with
the character $\chi:V_\pm\to k_0^\times$.
Indeed $D\subset k_0^\times$ and $\ell$ is $k_0$-linear,
which proves exactly the displayed multiplication
identity. All target spaces are finite and discrete,
so composition also preserves continuity. Equation
(F.25.8) gives $d_z(\mathcal K_\pm)\subset U$, and
therefore $\ell d_z$ vanishes on $\mathcal K_\pm$.
The character is nontrivial on $R_\pm$ because
$C\ne\{1\}$. Apply the proved F.24.1 restriction
injection over the finite field $k_0$. It gives
$\ell d_z=0$ on all of $V_\pm$, for every such $\ell$.
Together with (F.25.9) this proves
\[
 d_z(V_\pm)\subset U,\qquad
 N_z\subset U\rtimes C=M_z.                                  \tag{F.25.11}
\]
The inclusion $\mathcal K_\pm\leq R_\pm$ supplies the
reverse inclusion. Conjugating back by $\tau_z$ gives
the exact equality (F.25.2).

For the class $k_1=k_0$, the $k_0$-linearity of $U$
has already been proved using (F.25.6). Combining
this with the separate case $C=\{1\}$ proves
(F.25.3). For $k=\mathbb F_p$, both generated
subfields are $\mathbb F_p$ for every character,
so (F.25.4) follows for every prime, including two
and three. The proof applies to either of the actual
marginals, using its corresponding F.24 injection;
that injection already contains the explicit F.23
sign-symmetry comparison.

For the remaining class $k_1\subsetneq k_0$,
(F.25.6) supplies $k_1$-linearity of $U$. The
coefficient map (F.25.10) requires $k_0$-linearity.
The proof has established the passage precisely
for the representations specified in the statement,
and supplies no quotient of the actual ambient
module for an arbitrary such $U$.
\end{proof}
